\section{零售预测结论、适用边界与推广}

\subsection{四问结论}

问题一把最初按 $9800$ 条明细计订单的口径改为按订单号去重，确认实际订单数为 $4921$，并发现销售额最高一成明细贡献 $60.40\%$、第四季度贡献 $38.52\%$，销售呈现交易长尾与年末集中的双重结构。% src:交接/实验记录.json:[问题=问题1/订单口径].尝试,现象,决定;求解/问题1/结果/验证与灵敏度.json:朴素基线;求解/问题1/结果/核心统计与业务指标.json:明细层Sales分布.Top10%样本销售额占比;求解/问题1/结果/时间结构.json:季度季节结构

问题二最初把区域列入核心，完整参数网格复核后只保留产品历史价值层、产品大类和产品子类，三者均在 $81/81$ 组中入选；区域降为次级地理信号，只在 $54/81$ 组中入选。% src:交接/实验记录.json:[问题=问题2/正确性复核的稳定阈值灵敏度].尝试,现象,决定;求解/问题2/结果/关键参数灵敏度.json:各因素入选次数,组合总数;求解/问题2/结果/双证据筛选.json:稳定核心因素,阈值敏感的次级信号

问题三曾纳入低阶SARIMA候选，但有候选产生非有限预测且最终权重接近零，故部署由季节朴素、阻尼ETS与谐波岭构成的受约束组合；封存期MAPE为 $16.367\%$，低于季节朴素的 $24.865\%$。% src:交接/实验记录.json:[问题=问题3/复杂模型与区间表现].现象,决定;求解/问题3/结果/2025评估指标.json:各模型测试指标.锁定受约束组合.MAPE,各模型测试指标.季节朴素.MAPE

问题四放弃细层级独立建模后采用总量锚定的份额收缩，验证期平均份额WAPE由固定份额的 $39.585\%$ 降至 $29.015\%$；下一年度销售额点预测为 $791951.82$，全年可加总风险边界为 $[454803.67,1379541.60]$。2025年封存测试的 $95\%$ 经验区间覆盖率仅为 $83.33\%$，该边界不作严格概率保证。% src:交接/实验记录.json:[问题=问题4/层级预测方法选择].尝试,决定;求解/问题4/结果/份额收缩回测与灵敏度.json:2025五维平均对比;求解/问题4/结果/2026总量与一致性.json:2026全年总量预测;求解/问题3/结果/2025评估指标.json:区间测试表现.95%覆盖率;交接/结果解读_问题3.md:局限与使用边界

\subsection{适用边界}

这些结论适用于当前企业既有类别、相同订单口径和相近经营环境。数据缺少利润、成本、折扣、促销、销量与库存，产品和区域结果只能解释销售关联与容量优先级；可加总风险边界未联合刻画份额误差，细层级名次也不宜转化为刚性预算或因果承诺。2026年外推还以表\ref{tab:model-assumption-boundary}中的季节位置与份额机制延续假设为前提；出现结构突变时，须转入滚动重估。% src:交接/数据档案.json:列清单,给建模师的话;交接/结果解读_问题4.md:五.局限;交接/结果解读_问题3.md:局限与使用边界

\subsection{推广方向}

若后续接入价格、促销、库存、利润和门店覆盖数据，可在保持时间外切分与层级加总约束的前提下，扩展为需求弹性、利润目标和库存服务水平的联合决策；若出现新区域或新品类，则应先设置“其他/新层级”份额，再随新增观测滚动更新总量模型与收缩权重。% src:交接/计划.json:问题清单[编号=4].方法理由;交接/结果解读_问题4.md:五.局限
